\chapter{Literature Review}

\section{Introduction}

Private Set Intersection has been studied since the early 2000s, with protocols
evolving from public-key-based constructions to highly optimised Oblivious Transfer
(OT) extensions. This chapter surveys the most relevant lines of work: classical
OT-based PSI, homomorphic encryption-based PSI, laconic PSI from pairings, and
the lattice-based laconic PSI protocol that this project implements. We also review
the post-quantum cryptography standardisation landscape and identify the research
gaps that this project addresses.

Two important nuances regarding post-quantum status and prior laconic PSI work are
addressed carefully throughout this chapter, as they are frequently mischaracterised
in the broader PSI literature:

\begin{enumerate}
    \item \textbf{OT-based protocols such as KKRT16 are not unconditionally
          non-post-quantum.} Their base OT can be instantiated from lattice or
          code-based assumptions via OT extension, making the full pipeline
          potentially post-quantum. The fundamental limitation of these protocols
          relative to laconic PSI is $O(n + m)$ communication scaling, not
          quantum vulnerability per se.
    \item \textbf{Prior laconic PSI implementations exist.} Aranha, Lin, Orlandi,
          and Simkin (ALOS22)~\cite{alos2022} provide the first empirically evaluated
          laconic PSI protocol, constructed from pairings. This project is correctly
          positioned as the first \textit{post-quantum} laconic PSI implementation,
          complementing rather than superseding ALOS22.
\end{enumerate}

% ─────────────────────────────────────────────────────────────────────────────
\section{Classical PSI: OT-Based Protocols}

The most efficient PSI protocols in practice are based on Oblivious Transfer (OT)
and OT extension. These protocols achieve remarkable computational performance ---
processing millions of records in sub-second runtimes --- and represent the
practical state of the art for balanced datasets.

\subsection{KKRT16 --- Batched Oblivious PRF}
Kolesnikov, Kumaresan, Rosulek, and Trieu~\cite{kolesnikov2016kkrt} construct
PSI from batched Oblivious PRF (OPRF) atop OT extension, achieving $O(n + m)$
communication and sub-second runtimes for $n = m = 10^4$. KKRT is the primary
classical baseline in this project.

\textbf{Post-quantum status.} KKRT is typically instantiated with an
elliptic-curve-based base OT (e.g., Naor-Pinkas), which is vulnerable to Shor's
algorithm. However, OT extension amortises the cost of base OT across many instances,
meaning that replacing only the base OT with a lattice-based or code-based
alternative adds only modest overhead to the full pipeline, making a fully
post-quantum KKRT pipeline theoretically feasible. The fundamental limitation of
KKRT relative to laconic PSI is therefore the $O(n + m)$ communication complexity,
which grows linearly with the large server set size $m$ regardless of the base OT
instantiation. In our experimental comparison (Section~6.8), KKRT is evaluated in
its standard elliptic-curve instantiation, as no benchmarked post-quantum KKRT
implementation exists at comparable scale.

\subsection{OKVS-Based PSI}
Garimella et al.~\cite{garimella2023okvs} achieve 20--40\% improvement over KKRT
through Oblivious Key-Value Stores (OKVS), reducing ciphertext overhead while
maintaining $O(n + m)$ communication. Like KKRT, the base OT can in principle be
replaced with a post-quantum alternative. These protocols lack the sublinear
server message that makes laconic PSI attractive for highly unbalanced settings
where $m \gg n$.

\subsection{Circuit-PSI}
Pinkas et al.~\cite{pinkas2019circuit} combine garbled circuits with OPRF to
compute arbitrary functions over the intersection without revealing the intersection
itself. While powerful for computation-over-intersection scenarios, Circuit-PSI
incurs $O(nm)$ gate count and $O(n + m)$ communication, making it unsuitable for
the large, unbalanced datasets considered in this project.

\subsection{HE-Based PSI for Unbalanced Datasets}
Chen, Laine, and Rindal~\cite{chen2017fhe} apply Homomorphic Encryption
(specifically CKKS/BFV) to PSI in the unbalanced regime where $m \gg n$, achieving
sublinear communication in $n$. CKKS/BFV are Ring-LWE-based schemes and are
conjectured to be post-quantum secure --- Shor's algorithm provides no advantage
against them. However, unlike ML-KEM and ML-DSA, these FHE schemes have not
undergone NIST's formal PQC standardisation process, so their parameter sets have
received less public external security scrutiny. Furthermore, the server message in
these constructions still scales with $m$, whereas in laconic PSI the server's first
message is $O(\log m)$ and is reusable across all future client queries ---
a qualitatively different and stronger property.

% ─────────────────────────────────────────────────────────────────────────────
\section{Laconic PSI from Pairings (ALOS22)}

Aranha, Lin, Orlandi, and Simkin~\cite{alos2022} present the first empirically
validated laconic PSI protocol, constructed from pairing-friendly elliptic curves
(BLS12-381). This is the most directly comparable prior work to this project and
serves as the classical laconic PSI performance baseline.

\textbf{Construction.} ALOS22 builds on bilinear Diffie-Hellman (BDDH) and
$q$-SDH assumptions over BLS12-381 pairings. The server compresses its dataset
into a short digest using algebraic commitments, and the client encrypts each
query element under the server's public parameters. The server performs decryption
checks to identify intersection elements.

\textbf{Communication.} $O(\log m)$ --- fully laconic. The server's first message
grows only logarithmically with its dataset size, identical in asymptotic behaviour
to DKLLMR23.

\textbf{Performance.} ALOS22 reports millisecond-to-second runtimes for datasets
of millions of records, exploiting the algebraic efficiency of pairing operations
over BLS12-381. At $m = 10^4$, ALOS22 achieves approximately 0.1 minutes runtime,
compared to 153 minutes for our Ring-LWE implementation. This performance gap is
discussed in detail in Section~6.8.

\textbf{Post-quantum status.} ALOS22 is \textit{not} post-quantum secure.
Pairing-based constructions over elliptic curves are vulnerable to quantum Fourier
sampling --- a generalisation of Shor's algorithm applied to the elliptic curve
discrete logarithm problem --- meaning the BLS12-381 security assumption is broken
in polynomial time by a sufficiently powerful quantum computer.

\textbf{Significance for this project.} ALOS22 is the first paper to empirically
demonstrate that laconic PSI is practically viable, providing the classical
performance ceiling for laconic PSI. Our work is the direct post-quantum analogue:
we provide the first empirical benchmark of Ring-LWE-based laconic PSI, playing
the same foundational role for the post-quantum setting that ALOS22 played for
the classical setting.

\textbf{Performance gap analysis.} The approximately $1{,}000\times$ runtime overhead
of our implementation relative to ALOS22 decomposes as follows. Ring-LWE gadget
decomposition (GInvMNTT) expands each witness polynomial by a factor of
$\log_2 q \approx 58$, a cost entirely absent in pairing-based constructions
where witnesses are single group elements. This expansion is cryptographically
necessary for Ring-LWE security and cannot be reduced without replacing the
underlying construction. An additional $\approx 10\times$ overhead arises from
Merkle proof generation and witness fetching. This decomposition is quantified
empirically in Sections 6.6 and 6.8.

% ─────────────────────────────────────────────────────────────────────────────
\section{Laconic Encryption and PSI from Ring-LWE (DKLLMR23)}

D\"{o}ttling, Kolonelos, Lai, Lin, Malavolta, and Rahimi~\cite{dottling2023efficient}
construct Laconic Encryption (LE) from Ring-LWE and extend it to PSI. This is
the protocol implemented in this project.

\textbf{Construction overview.} The server builds a Merkle tree over its dataset
and generates a Ring-LWE key pair per leaf. The Merkle root and public keys form
a compact public parameter set (the digest). The client encrypts each of its
elements under the server's public parameters, producing a ciphertext that can
only be successfully decrypted if the encrypted element matches a server leaf.
The server performs decryption checks using witness vectors derived from the Merkle
tree. The formal definition of the laconic encryption primitive, including
\textsf{LaconicEnc} and its inputs, is given in Chapter 3.

\textbf{Hardness assumption.} Ring-LWE hardness --- widely believed to be post-quantum
secure and closely related to the Module-LWE assumption standardised in NIST FIPS 203
(ML-KEM)~\cite{nist2024mlkem}.

\textbf{Communication complexity.} $O(n \log m)$ --- sublinear in server dataset
size $m$. For $n = 10^3$ client records and $m = 10^6$ server records, this yields
approximately 3.4~MB versus 97.8~MB for KKRT --- a $29\times$ reduction, measured
empirically in Section 6.5.

\textbf{Computation complexity.} $O(m \log m)$ server setup; $O(nm \cdot D)$
intersection detection, where $D$ is the ring dimension.

\textbf{Critical gap.} The DKLLMR23 paper provides no implementation or empirical
benchmark. All performance claims are asymptotic. This project provides the first
complete implementation and empirical evaluation, establishing the first measured
performance data for Ring-LWE-based laconic PSI.

% ─────────────────────────────────────────────────────────────────────────────
\section{Post-Quantum Cryptography Standards}

NIST completed its Post-Quantum Cryptography standardisation process in 2024,
publishing:
\begin{itemize}
    \item \textbf{FIPS 203} (ML-KEM / Kyber)~\cite{nist2024mlkem}: Key encapsulation
          based on Module-LWE.
    \item \textbf{FIPS 204} (ML-DSA / Dilithium)~\cite{nist2024mldsa}: Digital
          signatures based on Module-LWE/SIS.
\end{itemize}

Both standards confirm the maturity and confidence of the research community in
lattice-based hardness assumptions for standardised cryptographic deployments. Our
implementation uses Ring-LWE --- the ring-structured variant of LWE --- which is
closely related to Module-LWE and shares the same underlying worst-case hardness
reduction to lattice problems~\cite{lyubashevsky2010ideal}. The NIST standardisation
process also considered isogeny-based schemes (SIKE), which were subsequently
broken by classical attacks in 2022, reinforcing the importance of lattice-based
approaches as the most robust post-quantum foundation currently available.

% ─────────────────────────────────────────────────────────────────────────────
\section{Research Gaps}

The survey above identifies the following specific gaps directly addressed by this project:

\begin{enumerate}
    \item \textbf{No open implementation of DKLLMR23.} Despite the theoretical
          construction existing since 2023, no open-source implementation exists
          and no empirical benchmark has been reported.
    \item \textbf{Memory scalability unknown.} The auxiliary space complexity
          $S_{\text{aux}}(m) = O(m \cdot D^2 \cdot \log q)$ makes memory management
          a non-trivial engineering problem at real-world scale. No prior work has
          characterised this experimentally or proposed a solution.
    \item \textbf{No post-quantum laconic PSI benchmark.} ALOS22 empirically
          validated classical laconic PSI. There is no prior empirical data on
          the performance cost of post-quantum laconic PSI, either at 64-bit or
          128-bit quantum security levels.
    \item \textbf{No real-world deployment study.} No prior work validates laconic
          PSI in a GDPR-compliance scenario with concrete infrastructure, API design,
          and audit trail requirements.
\end{enumerate}

% ─────────────────────────────────────────────────────────────────────────────
\section{Summary Comparison}

\begin{table}[H]
\centering
\caption{Comparison of PSI protocols}
\label{tab:psi-comparison}
\begin{tabular}{@{}llcccl@{}}
\toprule
\textbf{Protocol} & \textbf{Assumption} & \textbf{PQ?} & \textbf{Laconic?} & \textbf{Communication} & \textbf{Empirical?} \\
\midrule
KKRT16~\cite{kolesnikov2016kkrt}  & OT (base OT replaceable) & Possible$^\dagger$ & No  & $O(n + m)$ & Yes \\
OKVS~\cite{garimella2023okvs}     & OT (base OT replaceable) & Possible$^\dagger$ & No  & $O(n + m)$ & Yes \\
ALOS22~\cite{alos2022}            & Pairing (BLS12-381)      & No                 & Yes & $O(\log m)$ & Yes \\
HE-PSI~\cite{chen2017fhe}         & CKKS/BFV (lattice)       & Partial$^\ddagger$ & No  & $O(n)$      & Yes \\
DKLLMR23 (ours)                   & Ring-LWE                 & \textbf{Yes}       & \textbf{Yes} & \textbf{$O(n \log m)$} & \textbf{This work} \\
\bottomrule
\end{tabular}
\end{table}

\noindent{\small $^\dagger$ \textit{Post-quantum instantiation of OT-based protocols is theoretically feasible via lattice-based base OT with OT extension, but no benchmarked implementation at comparable scale exists.}}\\
\noindent{\small $^\ddagger$ \textit{CKKS/BFV are lattice-based but have not been standardised under NIST PQC; their post-quantum security margins are less precisely characterised than ML-KEM/ML-DSA.}}\\

Our protocol is the only entry in Table~\ref{tab:psi-comparison} that is
simultaneously \textbf{laconic} (sublinear server message, reusable across clients)
and \textbf{plausibly post-quantum secure} under a NIST-endorsed hardness assumption.

\clearpage
